\chapter{رشته‌ها و کاراکترها}%
\label{ch:strings}

متن همه‌جا هست: نام‌ها، پیام‌ها، خروجیِ برنامه‌ها. در این فصل عمیق‌تر با
\textbf{رشته}ها کار می‌کنیم: آن‌ها را به هم می‌چسبانیم، تکه‌تکه می‌کنیم،
می‌شماریم و دگرگون می‌کنیم.

\section{یادآوری: رشته چیست؟}

رشته دنباله‌ای از نویسه‌هاست که میانِ دو نقل‌قول نوشته می‌شود:
\code{"سلام دنیا"}. هر رشته از کنارِ‌هم‌قرارگرفتنِ نویسه‌ها (کاراکترها) ساخته
می‌شود. سلام از یونی‌کد پشتیبانی می‌کند، پس می‌توانید در رشته‌ها از فارسی،
انگلیسی، نشانه‌ها و حتی شکلک‌ها استفاده کنید.

\section{چسباندنِ رشته‌ها}

ساده‌ترین کار با رشته‌ها، چسباندنِ آن‌هاست. دو راه دارید: عملگرِ \code{+} یا
متدِ \code{.پیوست()}. این مثالِ واقعی هر دو راه را در عمل نشان می‌دهد:

\begin{salam}
// وارونه‌سازی رشته
تابع اصلی:
    س : رشته = "Salam"
    متغیر ر : رشته = ""
    متغیر i : صحیح = س.len() - 1
    تاوقتی i >= 0:
        ر = ر.concat(س.substr(i, 1))
        i = i - 1
    تمام
    چاپ "وارونه‌ی", س, "می‌شود", ر
تمام
\end{salam}

این برنامه یک رشته را \emph{وارونه} می‌کند. بیایید بشکافیمش:
\begin{itemize}
  \item \code{س.len()} طولِ رشته را می‌دهد؛ پس \code{س.len() - 1} اندیسِ آخرین
        نویسه است.
  \item حلقه از آخرین نویسه تا نخستین (\code{i >= 0}) عقب می‌رود.
  \item \code{س.substr(i, 1)} (معادلِ \code{.زیررشته(i, 1)}) یک
        \emph{زیررشته} به طولِ ۱ از موقعیتِ \code{i} می‌گیرد یعنی تنها
        یک نویسه.
  \item \code{ر.concat(...)} (معادلِ \code{.پیوست(...)}) آن نویسه را به انتهای
        \code{ر} می‌چسباند.
\end{itemize}
خروجی: \code{وارونه‌ی Salam می‌شود malaS}.

\begin{note}
متدهای رشته نیز مانندِ متدهای بردار هم نامِ فارسی دارند و هم انگلیسی:
\code{len}/\code{طول}، \code{substr}/\code{زیررشته}، \code{concat}/\code{پیوست}.
در یک برنامه بهتر است یک شیوه را برگزینید و به آن پایبند بمانید.
\end{note}

\section{پیمایشِ نویسه‌به‌نویسه}

بسیاری از کارهای متنی نیازمندِ آن‌اند که رشته را نویسه‌به‌نویسه بررسی کنیم. الگوی
آن همیشه یکسان است: یک حلقه از ۰ تا \code{طول - 1}، و در هر گام گرفتنِ یک
زیررشتهٔ تک‌نویسه‌ای. این مثالِ واقعی مصوت‌های یک متن را می‌شمارد:

\begin{salam}
// LANG: fa
تابع اصلی:
    س : رشته = "salam donya"
    متغیر تعداد : صحیح = 0
    تکرار س.طول() با i:
        ک : رشته = س.زیررشته(i, 1)
        اگر ک == "a" || ک == "e" || ک == "i" || ک == "o" || ک == "u":
            تعداد = تعداد + 1
        تمام
    تمام
    چاپ "تعداد مصوت‌ها در", س, "=", تعداد
تمام
\end{salam}

این الگو «برای هر نویسه، کاری انجام بده» را در شمارش، جست‌وجو، فیلتر
کردن و بسیاری کارهای دیگر به‌کار خواهید برد.

\begin{warn}
برای مقایسهٔ دو رشته از \code{==} استفاده کنید، نه از کارهای دیگر. در سلام
\code{"a" == "a"} درست است و دو رشته را نویسه‌به‌نویسه می‌سنجد. به یاد داشته
باشید که مقایسه به بزرگی و کوچکیِ حروف حساس است: \code{"A"} با \code{"a"}
برابر نیست.
\end{warn}

\section{ابزارهای آمادهٔ کتابخانهٔ رشته}

نوشتنِ همهٔ کارهای متنی از صفر خسته‌کننده است. خوشبختانه سلام یک کتابخانهٔ
آمادهٔ رشته دارد که با \code{واردکردن "str"} در دسترس قرار می‌گیرد. به این
مثالِ واقعی نگاه کنید:

\begin{salam}
// LANG: fa
واردکردن "str"

تابع اصلی:
    ثابت متن : رشته = "Salam Donya"
    چاپ "طول:", رشته.طول(متن)
    چاپ "بزرگ:", رشته.بزرگ(متن)
    چاپ "کوچک:", رشته.کوچک(متن)
    چاپ "وارونه:", رشته.وارونه("Salam")
    چاپ "پیوست:", رشته.پیوست("سلام ", "دنیا")
تمام
\end{salam}

با \code{واردکردن "str"} ماژولِ رشته را وارد می‌کنیم و سپس با نامِ فارسیِ آن،
\code{رشته}، به توابعش دسترسی پیدا می‌کنیم. مهم‌ترین توابع:

\begin{center}
\begin{tabular}{l l l}
\toprule
\textbf{فارسی} & \textbf{انگلیسی} & \textbf{کار} \\
\midrule
\code{رشته.طول(م)} & \code{str.Len} & طولِ رشته \\
\code{رشته.بزرگ(م)} & \code{str.Upper} & تبدیل به حروفِ بزرگ \\
\code{رشته.کوچک(م)} & \code{str.Lower} & تبدیل به حروفِ کوچک \\
\code{رشته.وارونه(م)} & \code{str.Reverse} & وارونه‌سازی \\
\code{رشته.پیوست(آ, ب)} & \code{str.Concat} & چسباندنِ دو رشته \\
\bottomrule
\end{tabular}
\end{center}

\begin{note}
دو شیوهٔ کار با رشته را دیدیم: یکی به‌صورتِ \emph{متد} (مثلِ \code{متن.طول()})
و دیگری به‌صورتِ \emph{تابعِ ماژول} (مثلِ \code{رشته.طول(متن)}). هر دو درست‌اند؛
ماژولِ \code{str} توابعِ پیشرفته‌تری (مانندِ جداسازی، جست‌وجو و تبدیل به عدد)
هم دارد که در فصلِ کتابخانهٔ استاندارد (فصلِ ۱۵) بیشتر می‌بینیم.
\end{note}

\section{دنباله‌های گریز}

گاهی می‌خواهیم نویسه‌هایی در رشته بگذاریم که نوشتنشان مستقیم دشوار است مثلاً
خودِ نقل‌قول، یا رفتن به خطِ تازه. برای این کار از \textbf{دنباله‌های گریز}
استفاده می‌کنیم که با یک \code{\textbackslash} (بک‌اسلش) آغاز می‌شوند:

\begin{center}
\begin{tabular}{c l}
\toprule
\textbf{دنباله} & \textbf{معنا} \\
\midrule
\lr{\texttt{\textbackslash n}} & خطِ تازه \\
\lr{\texttt{\textbackslash t}} & تب (فاصلهٔ افقی) \\
\lr{\texttt{\textbackslash"}} & نقل‌قولِ دوتایی \\
\lr{\texttt{\textbackslash\textbackslash}} & خودِ بک‌اسلش \\
\bottomrule
\end{tabular}
\end{center}

برای نمونه، \code{چاپ "خطِ یک\textbackslash nخطِ دو"} دو خط چاپ می‌کند، چون
\lr{\texttt{\textbackslash n}} میانِ آن‌ها یک شکستِ خط می‌گذارد.

\section{کاراکتر در برابرِ رشته}

در فصلِ ۲ دیدیم که کاراکتر (\code{'A'}) و رشته (\code{"A"}) دو چیزِ متفاوت‌اند.
کاراکتر تنها \emph{یک} نویسه است و میانِ تک‌نقل‌قول می‌آید، حال‌آنکه رشته
می‌تواند هیچ، یک، یا هزاران نویسه داشته باشد و میانِ نقل‌قولِ دوتایی می‌آید.
در عمل، برای کار با متن بیشتر با رشته‌ها سروکار دارید؛ کاراکتر بیشتر در کارهای
سطحِ‌پایین به کار می‌آید.

\section{مثالِ کامل: شمارشِ واژه‌ها}

بیایید آموخته‌ها را در یک برنامهٔ کوچک به کار ببندیم که نویسه‌های یک رشته را
می‌شمارد و تعدادِ فاصله‌ها را نیز می‌یابد که تقریباً برابرِ تعدادِ واژه‌هاست:

\begin{salam}
// LANG: fa
تابع اصلی:
    متن : رشته = "سلام دنیا و سلام بر شما"
    متغیر فاصله : صحیح = 0
    تکرار متن.طول() با i:
        اگر متن.زیررشته(i, 1) == " ":
            فاصله = فاصله + 1
        تمام
    تمام
    چاپ "تعداد نویسه‌ها:", متن.طول()
    چاپ "تعداد واژه‌ها:", فاصله + 1
تمام
\end{salam}

\section{خلاصهٔ فصل}

\begin{itemize}
  \item رشته دنباله‌ای از نویسه‌هاست و از یونی‌کد پشتیبانی می‌کند.
  \item با \code{+} یا \code{.پیوست()} رشته‌ها را به هم می‌چسبانیم.
  \item \code{.طول()} طول و \code{.زیررشته(i, n)} یک تکه از رشته را می‌دهد.
  \item الگوی پیمایشِ نویسه‌به‌نویسه برای شمارش و جست‌وجو بسیار کاربردی است.
  \item ماژولِ \code{str} توابعِ آماده‌ای مانندِ \code{بزرگ}، \code{کوچک} و
        \code{وارونه} دارد.
  \item دنباله‌های گریز مانندِ \lr{\texttt{\textbackslash n}} نویسه‌های ویژه
        را در رشته می‌گنجانند.
\end{itemize}

\section{تمرین‌ها}

\exercise برنامه‌ای بنویسید که یک رشته بگیرد و بگوید چند نویسه دارد.

\exercise با کمکِ ماژولِ \code{str}، یک نام را هم با حروفِ بزرگ و هم با حروفِ
کوچک چاپ کنید.

\exercise برنامه‌ای بنویسید که بشمارد حرفِ «ا» چند بار در یک متنِ فارسی آمده
است.

\exercise یک رشته را وارونه کنید (یک بار با حلقهٔ دستی و یک بار با
\code{رشته.وارونه}) و نتیجه‌ها را مقایسه کنید.

\exercise برنامه‌ای بنویسید که یک رشته را بگیرد و بررسی کند آیا «متقارن»
(palindrome) است یعنی از چپ و راست یکسان خوانده می‌شود.
